% ============================================================
%  Math Notes Template
% ============================================================
\documentclass[11pt, a4paper]{article}

% --- Mathematics (must load before newtxmath) ---
\usepackage{amsmath, amssymb, amsthm, mathtools}

% --- Fonts & Encoding (XeLaTeX) ---
\usepackage{fontspec}
\usepackage{libertine}                    % Linux Libertine text + Biolinum sans
\usepackage[libertine]{newtxmath}         % Matching math font (loads after amsthm)
\usepackage{inconsolata}                  % Monospace font
\usepackage{microtype}                    % Improved typography

% --- Page Layout ---
\usepackage[
  margin=2.5cm,
  headheight=14pt
]{geometry}

% --- Colours ---
\usepackage[dvipsnames]{xcolor}
\definecolor{defblue}{HTML}{1a5276}
\definecolor{thmgreen}{HTML}{196f3d}
\definecolor{exred}{HTML}{922b21}
\definecolor{linkblue}{HTML}{2471a3}

% --- TikZ ---
\usepackage{tikz}
\usetikzlibrary{
  arrows.meta,
  calc,
  decorations.markings,
  patterns,
  positioning,
  shapes.geometric,
  cd                                      % commutative diagrams
}

% --- Theorem Environments ---
\usepackage{thmtools}

\declaretheoremstyle[
  headfont=\bfseries\color{thmgreen},
  bodyfont=\normalfont,
  spaceabove=10pt,
  spacebelow=10pt,
]{thmstyle}

\declaretheoremstyle[
  headfont=\bfseries\color{defblue},
  bodyfont=\normalfont,
  spaceabove=10pt,
  spacebelow=10pt,
]{defstyle}

\declaretheoremstyle[
  headfont=\bfseries\color{exred},
  bodyfont=\normalfont,
  spaceabove=10pt,
  spacebelow=10pt,
]{exstyle}

\declaretheorem[style=thmstyle, numberwithin=section, name=Theorem]{theorem}
\declaretheorem[style=thmstyle, sibling=theorem, name=Lemma]{lemma}
\declaretheorem[style=thmstyle, sibling=theorem, name=Proposition]{proposition}
\declaretheorem[style=thmstyle, sibling=theorem, name=Corollary]{corollary}
\declaretheorem[style=defstyle, sibling=theorem, name=Definition]{definition}
\declaretheorem[style=exstyle,  sibling=theorem, name=Example]{example}
\declaretheorem[style=exstyle,  numbered=no,      name=Remark]{remark}

% --- Hyperlinks & Cross-references ---
\usepackage[
  colorlinks=true,
  linkcolor=linkblue,
  citecolor=linkblue,
  urlcolor=linkblue
]{hyperref}
\usepackage{cleveref}

% --- Headers & Footers ---
\usepackage{fancyhdr}
\pagestyle{fancy}
\fancyhf{}
\fancyhead[L]{\textsc{\nouppercase{\leftmark}}}
\fancyhead[R]{\thepage}
\renewcommand{\headrulewidth}{0.4pt}

% --- Enumeration ---
\usepackage{enumitem}

% --- Bibliography (optional, uncomment if needed) ---
% \usepackage[style=alphabetic, backend=biber]{biblatex}
% \addbibresource{refs.bib}

% ============================================================
%  Fill these in for each set of notes
% ============================================================
\newcommand{\notestitle}{Sheaf Theory}
\newcommand{\notesauthor}{B. Arm}
\newcommand{\notesdate}{\today}

\title{\notestitle}
\author{\notesauthor}
\date{\notesdate}

% ============================================================
\begin{document}

\maketitle
\tableofcontents
\newpage

% ============================================================
%  Chapter I: Sheaves and Presheaves
% ============================================================
\section{Sheaves and Presheaves}

\subsection{Definitions}

\subsection{Homomorphisms, Subsheaves, and Quotient Sheaves}

\subsection{Direct and Inverse Images}

\subsection{Cohomomorphisms}

\subsection{Algebraic Constructions}

\subsection{Supports}

\subsection{Classical Cohomology Theories}

% ============================================================
%  Chapter II: Sheaf Cohomology
% ============================================================
\section{Sheaf Cohomology}

\subsection{Differential Sheaves and Resolutions}

\subsection{The Canonical Resolution and Sheaf Cohomology}

\subsection{Injective Sheaves}

\subsection{Acyclic Sheaves}

\subsection{Flabby Sheaves}

\subsection{Connected Sequences of Functors}

\subsection{Axioms for Cohomology and the Cup Product}

\subsection{Maps of Spaces}

\subsection{$\Phi$-Soft and $\Phi$-Fine Sheaves}

\subsection{Subspaces}

\subsection{The Vietoris Mapping Theorem and Homotopy Invariance}

\subsection{Relative Cohomology}

\subsection{Mayer--Vietoris Theorems}

\subsection{Continuity}

\subsection{The K\"unneth and Universal Coefficient Theorems}

\subsection{Dimension}

\subsection{Local Connectivity}

\subsection{Change of Supports; Local Cohomology Groups}

\subsection{The Transfer Homomorphism and the Smith Sequences}

% ============================================================
%  Chapter III: Comparison with Other Cohomology Theories
% ============================================================
\section{Comparison with Other Cohomology Theories}

\subsection{Singular Cohomology}

\subsection{Alexander--Spanier Cohomology}

\subsection{de Rham Cohomology}

\subsection{\v{C}ech Cohomology}

% ============================================================
%  Chapter IV: Applications of Spectral Sequences
% ============================================================
\section{Applications of Spectral Sequences}

\subsection{The Spectral Sequence of a Differential Sheaf}

\subsection{The Fundamental Theorems of Sheaves}

\subsection{Direct Images}

\subsection{The Leray Sheaf}

\subsection{Extension of a Support Family by a Family on the Base Space}

\subsection{The Leray Spectral Sequence of a Map}

\subsection{Fiber Bundles}

\subsection{Dimension}

\subsection{The Fary Spectral Sequence}

\subsection{Sphere Bundles with Singularities}

\subsection{The Oliver Transfer and the Conner Conjecture}

% ============================================================
%  Chapter V: Borel--Moore Homology
% ============================================================
\section{Borel--Moore Homology}

\subsection{Cosheaves}

\subsection{The Dual of a Differential Cosheaf}

\subsection{Homology Theory}

\subsection{Maps of Spaces}

\subsection{Subspaces and Relative Homology}

\subsection{The Vietoris Theorem, Homotopy, and Covering Spaces}

\subsection{The Homology Sheaf of a Map}

\subsection{The Basic Spectral Sequences}

\subsection{Poincar\'e Duality}

\subsection{The Cap Product}

\subsection{Intersection Theory}

\subsection{Uniqueness Theorems}

\subsection{Uniqueness Theorems for Maps and Relative Homology}

\subsection{The K\"unneth Formula}

\subsection{Change of Rings}

\subsection{Generalized Manifolds}

\subsection{Locally Homogeneous Spaces}

\subsection{Homological Fibrations and $p$-adic Transformation Groups}

\subsection{The Transfer Homomorphism in Homology}

% ============================================================
%  Chapter VI: Cosheaves and \v{C}ech Homology
% ============================================================
\section{Cosheaves and \v{C}ech Homology}

\subsection{Theory of Cosheaves}

\subsection{Local Triviality}

\subsection{Local Isomorphisms}

\subsection{\v{C}ech Homology}

\subsection{The Reflector}

\subsection{Spectral Sequences}

% \printbibliography

\end{document}
